% §3 Two Encoding Constructions
%
% Merges old §3 (index-set) and §4 (tree-encoding) into a single
% side-by-side presentation.  Fixes the formula inconsistency between
% eq (7) (SRL) and the code/proof formula.

We now define the two parameterized maps that generate fermion-to-qubit
encodings.  Both factor through the Majorana intermediate representation
(\cref{def:majorana}), separating encoding \emph{definition} from
encoding \emph{computation}.

%%====================================================================
\subsection{Construction~A: Index-set schemes}
\label{sec:index-set}

\begin{definition}[Index-set scheme]
\label{def:index-set-scheme}
An \emph{index-set encoding scheme} for $n$ modes is a triple of
set-valued functions $\scheme = (U, P, \mathrm{Occ})$:
\begin{align}
  U &: \{0, \ldots, n{-}1\} \to 2^{\{0,\ldots,n{-}1\}},
    \label{eq:update-fn} \\
  P &: \{0, \ldots, n{-}1\} \to 2^{\{0,\ldots,n{-}1\}},
    \label{eq:parity-fn} \\
  \mathrm{Occ} &: \{0, \ldots, n{-}1\} \to 2^{\{0,\ldots,n{-}1\}}.
    \label{eq:occ-fn}
\end{align}
We call $U(j)$ the \emph{update set}, $P(j)$ the \emph{parity set},
and $\mathrm{Occ}(j)$ the \emph{occupation set} of mode~$j$.
\end{definition}

\paragraph{The encoding algorithm.}
Given a scheme~$\scheme$, the Majorana operators for mode~$j$ are:
\begin{align}
  c_j &= X_{\{j\} \cup U(j)} \;\cdot\; Z_{P(j)},
  \label{eq:cj-from-scheme} \\[4pt]
  d_j &= Y_j \;\cdot\; X_{U(j)} \;\cdot\;
    Z_{(P(j)\,\triangle\,\mathrm{Occ}(j))\,\setminus\,\{j\}},
  \label{eq:dj-from-scheme}
\end{align}
where $X_S$ denotes $X$ on every qubit in set~$S$ (identity elsewhere),
$\triangle$ is the symmetric difference, and qubit positions in
multiple sets are resolved left-to-right (the first assignment
wins).  The sets $\{j\} \cup U(j)$ and $P(j)$ are required to be
disjoint for the formula to be unambiguous; we verify this as part of
the correctness check.

The ladder operators are then given by~\eqref{eq:encoded-ladder}.

\begin{remark}[Relation to the SRL formula]
\label{rem:srl}
Seeley, Richard, and Love~\cite{seeley2012} present a more complex
formula that places $Y$ at positions in $\mathrm{Occ}(j) \cap P(j)$
and $X$ at $\mathrm{Occ}(j) \setminus P(j)$.  For the three classical
schemes (JW, BK, Parity), the sets $\{j\} \cup U(j)$ and $P(j)$ are
disjoint, and both formulas produce identical Majorana operators.  Our
simpler formulation~\eqref{eq:cj-from-scheme}--\eqref{eq:dj-from-scheme}
suffices for the structural analysis of \cref{sec:star-tree} and
matches the implementation.
\end{remark}

\paragraph{Three classical encodings as set triples.}
\begin{align}
  \text{JW:}\;&\;
    U(j) = \emptyset,\;
    P(j) = \{0,\ldots,j{-}1\},\;
    \mathrm{Occ}(j) = \{j\}.
    \label{eq:jw-scheme} \\[3pt]
  \text{Parity:}\;&\;
    U(j) = \{j{+}1,\ldots,n{-}1\},\;
    P(j) = \{j{-}1\}^*,\;
    \mathrm{Occ}(j) = \{j{-}1,j\}^*.
    \label{eq:parity-scheme} \\[3pt]
  \text{BK:}\;&\;
    U, P, \mathrm{Occ}\text{ from the Fenwick tree}
    ~\cite{fenwick1994, seeley2012}.
    \label{eq:bk-scheme}
\end{align}
(Starred entries: for $j = 0$, $P(0) = \emptyset$ and
$\mathrm{Occ}(0) = \{0\}$.)

All variation between encodings is captured in the three
functions.  The encoding algorithm itself never changes.

\paragraph{Deriving index sets from a tree.}
Given a labelled rooted tree $T$ on $[n]$, one can attempt to derive
a scheme by reading off structural relationships:
\begin{align}
  U(j) &= \text{ancestors of } j, \notag \\
  P(j) &= \text{remainder}(j) \cup \text{children}(j), \notag \\
  \mathrm{Occ}(j) &= \text{descendants}(j) \cup \{j\},
  \label{eq:tree-derived-sets}
\end{align}
where $\text{remainder}(j)$ collects children of ancestors of $j$
that have index $< j$ and are not themselves ancestors of~$j$.

This derivation always produces well-formed sets and yields Pauli
strings of the correct width.  The question is whether the resulting
encoding satisfies the CAR.  We prove in \cref{sec:star-tree} that it
does \emph{if and only if} the tree is a star.

%%====================================================================
\subsection{Construction~B: Path-based tree encodings}
\label{sec:tree-encoding}

\begin{definition}[Encoding tree]
\label{def:encoding-tree}
An \emph{encoding tree} for $n$ fermionic modes is a rooted tree $T$
with internal nodes, where each node has at most three descending
connections, each labelled by a distinct Pauli operator from $\{X, Y,
Z\}$.  Each connection is either an \emph{edge} (to a child node) or a
\emph{leg} (a dangling link).  The total number of legs is $2n$, paired
into $n$ mode assignments.
\end{definition}

\paragraph{The encoding algorithm.}
For mode~$j$, let $\ell_j^{(x)}$ and $\ell_j^{(y)}$ be the paired
legs.  The Majorana operator for each leg is the Pauli string obtained
by collecting the link labels along the root-to-leg path:
\begin{equation}
  S(\ell) = \bigotimes_{v \in \mathrm{path}(\ell)}
    \text{label}(v \to \text{next}) \;\otimes\;
    \bigotimes_{v \notin \mathrm{path}(\ell)} I.
  \label{eq:path-pauli}
\end{equation}
The ladder operators are:
$\ad{j} = \tfrac{1}{2}(S(\ell_j^{(x)}) - i\,S(\ell_j^{(y)}))$.

\begin{theorem}[Tree encoding correctness]
\label{thm:tree-car}
For any encoding tree $T$ with $2n$ legs paired into $n$ modes, the
ladder operators satisfy the CAR.
\end{theorem}

\begin{proof}[Proof sketch]
Legs on distinct subtrees produce Pauli strings with disjoint
non-identity support except at the branching ancestor, where distinct
labels ($X$, $Y$, $Z$) ensure anticommutation.  Legs sharing a prefix
differ in the subtree below the branch point.  The proof follows Jiang
et al.~\cite{jiang2020} by induction on tree depth.
\end{proof}

\paragraph{Recovery of known encodings.}
All five standard encodings arise from specific tree topologies:
\begin{center}
\begin{tabular}{@{}lll@{}}
\toprule
Tree topology & Encoding & Max Pauli weight \\
\midrule
Star (depth 1, chain ordering) & Jordan--Wigner & $O(n)$ \\
Star (depth 1, reversed) & Parity & $O(n)$ \\
Fenwick tree & Bravyi--Kitaev & $O(\log_2 n)$ \\
Balanced binary tree & Binary tree & $O(\log_2 n)$ \\
Balanced ternary tree & Ternary tree & $O(\log_3 n)$ \\
\bottomrule
\end{tabular}
\end{center}

%%====================================================================
\subsection{Pauli-weight bounds}
\label{sec:weight-bounds}

\begin{theorem}[Weight bound]
\label{thm:weight-bound}
For an encoding tree $T$ with depth $d = \depth{T}$, every encoded
ladder operator has Pauli weight at most $2d + 1$.
\end{theorem}

\begin{proof}
Each Majorana string has weight equal to the root-to-leg path length,
at most $d + 1$.  The worst-case ladder operator (sum of two Majorana
strings sharing only the root) has weight at most $2d + 1$.
\end{proof}

\begin{corollary}[Optimal asymptotic weight]
\label{cor:optimal-weight}
Since each node has at most 3 children (one per Pauli label), the
balanced ternary tree achieves optimal weight $O(\log_3 n)$ among
fixed-arity trees.
\end{corollary}

%%====================================================================
\subsection{Three constructions, not two}
\label{sec:three-constructions}

A key observation is that the literature contains three distinct
encoding constructions, each with a different domain of validity:

\begin{description}
  \item[Construction~A] (index-set method,
    \cref{sec:index-set}): derives index sets from a tree via
    \eqref{eq:tree-derived-sets} and applies the universal
    algorithm~\eqref{eq:cj-from-scheme}--\eqref{eq:dj-from-scheme}.
    Valid only for star trees (\cref{sec:star-tree}).

  \item[Construction~B] (path-based method,
    \cref{sec:tree-encoding}): generates Pauli strings directly from
    root-to-leg paths.  Universal: valid for every encoding tree.

  \item[Construction~F] (Fenwick-specific): the Bravyi--Kitaev
    encoding uses hand-derived bit-manipulation formulas for $U$, $P$,
    $\mathrm{Occ}$ that exploit the unique algebraic structure of the
    Fenwick tree~\cite{fenwick1994}.  These formulas produce a correct
    encoding but are not the same as the generic tree-derived sets of
    Construction~A applied to a Fenwick-shaped tree.
\end{description}

The SRL ``unification'' of JW, BK, and Parity into a single index-set
framework is, on inspection, a conflation of Constructions~A and~F.
JW and Parity are genuine Construction~A instances (star trees).  BK
is a Construction~F instance that happens to be \emph{expressible} as
an index-set scheme, but the scheme is hand-derived for the Fenwick
tree rather than mechanically extracted.
